% !TeX root = ../../main.tex

\section{Dominios de aplicación de Big Data}

\subsection{Introducción}

Big Data permite almacenar y analizar grandes cantidades de datos provenientes de diferentes fuentes. Su aplicación ayuda a descubrir patrones, hacer predicciones y tomar mejores decisiones en distintos sectores.

\subsection{Web}

En el dominio web, Big Data se utiliza para analizar la actividad de los usuarios en páginas y aplicaciones. Algunos ejemplos son:

\begin{itemize}
    \item Medir las visitas y el tiempo que los usuarios permanecen en un sitio.
    \item Evaluar el rendimiento de páginas y aplicaciones.
    \item Mostrar publicidad relacionada con los intereses del usuario.
    \item Recomendar películas, música, videos o noticias.
\end{itemize}

\subsection{Finanzas}

Los bancos y las instituciones financieras analizan datos de cuentas, pagos y movimientos para reducir riesgos. Sus principales aplicaciones son:

\begin{itemize}
    \item Calcular el riesgo de otorgar un crédito.
    \item Detectar operaciones con tarjetas de crédito que pueden ser fraudulentas.
    \item Identificar lavado de dinero.
    \item Analizar reclamaciones de seguros.
\end{itemize}

\subsection{Salud}

En el sector salud, Big Data permite combinar expedientes clínicos, resultados de laboratorio y datos de dispositivos médicos. Se utiliza para:

\begin{itemize}
    \item Detectar y vigilar brotes de enfermedades.
    \item Comparar pacientes con características similares.
    \item Predecir reacciones negativas a medicamentos.
    \item Detectar anomalías en reclamaciones de seguros médicos.
    \item Monitorear pacientes mediante dispositivos portátiles.
\end{itemize}

\subsection{Internet de las cosas}

El Internet de las cosas está formado por objetos y sensores conectados a Internet. Estos dispositivos producen datos que pueden analizarse mediante Big Data. Algunas aplicaciones son:

\begin{itemize}
    \item Detectar intrusiones mediante cámaras y sensores.
    \item Identificar espacios disponibles en estacionamientos.
    \item Monitorear el tráfico y las condiciones de las carreteras.
    \item Detectar daños en puentes y edificios.
    \item Controlar sistemas de riego de acuerdo con la humedad del suelo.
\end{itemize}

\subsection{Medio ambiente}

Los sistemas ambientales producen datos mediante sensores instalados en diferentes lugares. Big Data permite analizarlos para:

\begin{itemize}
    \item Monitorear las condiciones del clima.
    \item Medir la contaminación del aire y el ruido.
    \item Detectar incendios forestales.
    \item Generar alertas por inundaciones.
    \item Vigilar la calidad del agua.
\end{itemize}

\subsection{Logística y transporte}

En logística y transporte, Big Data ayuda a conocer la ubicación y condición de vehículos y mercancías. Sus aplicaciones incluyen:

\begin{itemize}
    \item Rastrear vehículos mediante GPS.
    \item Vigilar la temperatura y humedad de los productos transportados.
    \item Detectar fallas mecánicas en los vehículos.
    \item Calcular rutas más rápidas.
    \item Asignar conductores o repartidores cercanos.
    \item Ajustar los precios de taxis según la demanda y el tráfico.
\end{itemize}

\subsection{Industria}

En la industria, los sensores instalados en las máquinas producen datos sobre temperatura, vibración y funcionamiento. Esta información se utiliza para:

\begin{itemize}
    \item Predecir fallas antes de que ocurran.
    \item Determinar la causa de una avería.
    \item Detectar condiciones peligrosas para los trabajadores.
    \item Mejorar la planeación y el control de la producción.
\end{itemize}

\subsection{Comercio minorista}

Las tiendas utilizan Big Data para aumentar sus ventas y mejorar la atención al cliente. Algunas aplicaciones son:

\begin{itemize}
    \item Controlar el inventario mediante etiquetas RFID.
    \item Recomendar productos según las compras anteriores.
    \item Ofrecer promociones personalizadas.
    \item Organizar los productos dentro de la tienda.
    \item Pronosticar la demanda y evitar el exceso o la falta de productos.
\end{itemize}

\subsection{Conclusión}

Big Data tiene aplicaciones en numerosos sectores. En general, permite utilizar grandes cantidades de información para mejorar los servicios, reducir riesgos, aumentar la eficiencia y tomar decisiones con mayor rapidez.